\documentclass[10pt,a4paper]{amsart}
\usepackage[margin=19mm]{geometry}
\usepackage{amsmath,amssymb,mathtools}
\usepackage[scr=rsfs,cal=euler]{mathalfa}
\usepackage{xfrac}
\usepackage[unicode,hidelinks,bookmarks=false]{hyperref}
\newcommand{\ie}{i.e.\ }
\newcommand{\cf}{cf.\ }
\newcommand{\resp}{respectively\ }
\newcommand{\Subsection}{\subsection}
\providecommand{\nobreakdash}{\nobreak\hbox{-}}
\setlength{\emergencystretch}{2em}
\multlinegap0pt
\usepackage{amssymb}
\usepackage{enumerate}
\usepackage{braket}
\newtheorem{lem}{Lemma}
\newcommand{\End}{\operatorname{End}}
\newcommand{\adj}{\textrm{adj}}
\newcommand{\trm}[1]{\textrm{#1}}
\title{Additive Decompositions by Conjugacy Classes in $M_n(\mathbb{F}_q)$}
\author{Krishna Kishore, Sunil Kumar Mallick}
\date{Source: arXiv:2608.20736v1 (CC BY 4.0)}
\begin{document}
\maketitle

\noindent\emph{Source and licence.} Additive Decompositions by Conjugacy Classes in $M_n(\mathbb{F}_q)$. Version arXiv:2608.20736v1 (CC BY 4.0), distributed by arXiv under the Creative Commons Attribution 4.0 International licence, http://creativecommons.org/licenses/by/4.0/, as recorded in the arXiv licence metadata for this exact version and shown on the abstract page https://arxiv.org/abs/2608.20736v1. Full licence notice: \texttt{licenses/arxiv-2608.20736-CC-BY-4.0.txt}.\par\smallskip

\noindent\textit{Editorial scope.} Counted unit: the article's lem cyclic\_vector\_lemma (source0.tex lines 659--666). The article's complete original proof of it is delivered with it (source0.tex lines 667--697, one \texttt{proof} environment of 31 lines). No mathematics is written, repaired or re-derived: the statement and the proof are the article's own lines, in the article's own notation, and no retained line is substituted or re-flowed.  No further source-local context is needed: every cross-reference of every retained line resolves inside the two delivered spans.  Nothing was omitted from any retained span, so the package carries no omission record.  No retained line contains a citation, so the wrapper carries no bibliography and no bibliography entry.  No retained line contains a figure environment, an image-inclusion command or drawing code, so no image is delivered and the package declares no asset dependency.  Editorial formatting only, in addition: an AMS \texttt{amsart} preamble that restates the article's own declarations and macros used by the retained lines, namely the environments \texttt{lem} on separate counters, together with the standard packages those lines need.  The retained environments print in this document as Lemma 1 in order of appearance, whereas the article prints the same items with the numbering of its own full document.  The complete native source of the article as distributed in the arXiv e-print for this exact version is bundled unchanged as \texttt{sources/208204/source0.tex}.
\begin{lem}\label{cyclic_vector_lemma}
Let $W$ be a subspace of $V$ of dimension $m$, and $S \in \End(W)$, and $c \in W$. Define
$$
\Phi_{S,c} : W^* \to F[T]_{\leq m-1}
$$
by $\Phi_{S,c}(r) = r(\adj(TI-S)c)$; since the degree of $\adj(TI-S)c$ is at most $m-1$, this map is well-defined.
Then $\Phi_{S,c}$ is an isomorphism if and only if $c$ is a cyclic vector for $S$, i.e. $\Set{c, Sc, S^2 c, \ldots, S^{m-1}c}$ forms a basis of $W$.
\end{lem}

\begin{proof}
Let 
$$
Q(T) := \det(T I -S) = T^m +q_1 T^{m-1} + q_2 T^{m-2} + \ldots + q_m.
$$
Recall the standard identity 
\begin{align*}
\trm{adj}(TI-S) &= T^{m-1}I + T^{m-2}(S + q_1 I) + T^{m-3}(S^2 + q_1 S + q_2I) + \ldots + \\
& (S^{m-1} + q_1 S^{m-2} + \ldots + q_{m-1}I).
\end{align*}
Then 
$$
\trm{adj}(TI-S)c = T^{m-1}v_0 + T^{m-2}v_1 + \ldots + v_{m-1},
$$
where $v_0 = c$, $v_1 = Sc + q_1 c$, $v_2 = S^2 c + q_1 Sc + q_2 c$, and in general 
$$
v_j = S^j c + \trm{ a  linear combination of } c, Sc, \ldots, S^{j-1}c.
$$
Since $v_0, v_1,\ldots, v_{m-1}$ is obtained from the list $c, Sc, \ldots, S^{m-1}c$ by an upper triangular change of basis with diagonal entries $1$, it follows that $v_0, \ldots, v_{m-1}$ form a basis of $W$ if and only if $c, Sc, \ldots, S^{m-1}c$ form a basis of $W$.

Now 
\begin{align*}
\Phi_{S,c}(r) &= r ( \adj(TI-S)c ) = r \left( T^{m-1}v_0 + \ldots + v_{m-1} \right) \\
&= r(v_0) T^{m-1} + r(v_1) T^{m-2} + \ldots + r(v_{m-1}).
\end{align*}
The map 
$$
r \mapsto (r(v_0), \ldots, r(v_{m-1}))
$$
is an isomorphism of $W^*$ to $F^m$ if and only if $v_0, \ldots, v_{m-1}$ form a basis of $W$. Indeed, if the map is an isomorphism and $\sum_{i = 0}^{m-1} a_i v_i = 0$ then choosing $r_i$ such that $r_i(v_j) = \delta_{ij}$ we see that each $a_i = 0$ , hence $v_0, \ldots, v_{m-1}$ are linearly independent, and since they are $m$ in number, they form a basis of $W$. Conversely, if $v_0, \ldots, v_{m-1}$ form a basis of $W$, then choosing a basis of $W^*$ dual to $v_0, \ldots, v_{m-1}$ we see that the standard basis vectors of $F^m$ are contained in the image, so that the map is an isomorphism. The result follows.
\end{proof}

\end{document}
